\documentclass[11pt]{article}
\usepackage[margin=1.1in]{geometry}
\usepackage{graphicx}
\usepackage{booktabs}
\usepackage{longtable}
\usepackage{amsmath,amssymb,amsfonts}
\usepackage[section]{placeins}
\usepackage{url}
\usepackage[colorlinks=true,linkcolor=blue,citecolor=blue]{hyperref}

\newcommand{\botrule}{\bottomrule}
\renewcommand{\thetable}{S\arabic{table}}
\renewcommand{\thefigure}{S\arabic{figure}}
\renewcommand{\thesection}{S\arabic{section}}

\title{Supplementary Material\\
\large Learned initialization for variational quantum circuits on lattice and field-theory Hamiltonians: a ten-system benchmark of Q-MAML}
\author{Arnav Kapoor}
\date{}

\begin{document}
\maketitle

This document collects the per-seed numbers behind the means reported in the main text, and three
exploratory results that motivated choices made there but are not needed to follow its argument.

\section{Per-seed results for the ten Hamiltonians}\label{sec:supp_tracks}

Table~\ref{tab:supp_tracks} gives the raw per-seed relative error (or log-gap, Track F) underlying
Table~2 of the main text: every number that mean and win/tie/loss counts there are computed from.

\begin{longtable}{@{}llcccccc@{}}
\caption{Per-seed relative error (\%; log-gap for Track F) for every scheme and system at base size.}\label{tab:supp_tracks}\\
\toprule
Track & System & Seed & Q-MAML & Zero & $\pi$ & Uniform & Gaussian \\
\midrule
\endfirsthead
\toprule
Track & System & Seed & Q-MAML & Zero & $\pi$ & Uniform & Gaussian \\
\midrule
\endhead
A & Schwinger, $L{=}2$ & 0 & $2.0\!\times\!10^{-3}$ & 396 & 0.056 & 1.61 & 0.194 \\
 &  & 1 & 0.277 & 245 & 0.288 & 2.14 & 0.529 \\
 &  & 2 & 0.264 & 292 & 0.304 & 1.40 & 0.636 \\
B & $\mathbb{Z}_2$ LGT, penalty ansatz & 0 & 59.16 & 112 & 24.69 & 23.07 & 24.05 \\
 &  & 1 & 30.01 & 110 & 31.88 & 15.18 & 21.37 \\
 &  & 2 & 58.96 & 113 & 25.47 & 22.76 & 12.88 \\
B & $\mathbb{Z}_2$ LGT, gauge-inv.\ ansatz & 0 & $3.1\!\times\!10^{-3}$ & 0.016 & 0.027 & 0.022 & 0.015 \\
 &  & 1 & $2.9\!\times\!10^{-3}$ & $9.6\!\times\!10^{-3}$ & $8.8\!\times\!10^{-3}$ & $9.0\!\times\!10^{-3}$ & 0.010 \\
 &  & 2 & $6.6\!\times\!10^{-4}$ & $8.5\!\times\!10^{-3}$ & $7.5\!\times\!10^{-3}$ & $7.3\!\times\!10^{-3}$ & 0.012 \\
C & $\phi^4$ scalar field & 0 & 5.86 & 5.77 & 36.06 & 7.11 & 7.55 \\
 &  & 1 & 4.19 & 4.42 & 14.58 & 4.60 & 4.74 \\
 &  & 2 & 8.11 & 8.16 & 30.38 & 9.00 & 8.26 \\
D & SU(2) LGT, $N{=}2$ & 0 & $2.1\!\times\!10^{-5}$ & 1390 & 0.020 & 0.038 & 0.023 \\
 &  & 1 & $9.5\!\times\!10^{-6}$ & 775 & 0.011 & $8.5\!\times\!10^{-3}$ & 0.012 \\
 &  & 2 & $9.7\!\times\!10^{-6}$ & 1767 & 0.015 & 0.019 & $5.9\!\times\!10^{-3}$ \\
E & Neutrinos, $N{=}4$ & 0 & $2.2\!\times\!10^{-5}$ & 0.016 & $4.1\!\times\!10^{-4}$ & $4.4\!\times\!10^{-3}$ & $2.7\!\times\!10^{-3}$ \\
 &  & 1 & 0.011 & 0.078 & $5.1\!\times\!10^{-3}$ & 0.096 & 0.051 \\
 &  & 2 & $8.0\!\times\!10^{-5}$ & 0.073 & $7.2\!\times\!10^{-4}$ & $7.8\!\times\!10^{-4}$ & $4.1\!\times\!10^{-3}$ \\
F & SUSY QM (log-gap) & 0 & -4.87 & 0.00 & -2.40 & -4.75 & -4.66 \\
 &  & 1 & -4.90 & -0.58 & -3.28 & -4.62 & -4.62 \\
 &  & 2 & -3.67 & -0.05 & -1.95 & -4.89 & -4.35 \\
G & $\mathbb{D}_8$ LGT & 0 & 10.49 & 95.27 & 26.01 & 11.86 & 10.97 \\
 &  & 1 & 15.47 & 92.37 & 21.99 & 13.94 & 19.60 \\
 &  & 2 & 6.47 & 94.04 & 16.31 & 6.35 & 5.67 \\
H & Light nuclei (UCCSD) & 0 & 0.885 & 0.885 & 0.885 & 0.885 & 0.885 \\
 &  & 1 & 1.41 & 1.41 & 1.41 & 1.41 & 1.41 \\
 &  & 2 & 2.02 & 2.02 & 2.02 & 2.02 & 2.02 \\
I & Scalar Yukawa & 0 & $2.7\!\times\!10^{-5}$ & $3.0\!\times\!10^{-5}$ & 1.37 & $3.0\!\times\!10^{-4}$ & $3.7\!\times\!10^{-3}$ \\
 &  & 1 & $2.7\!\times\!10^{-5}$ & $2.9\!\times\!10^{-5}$ & 0.550 & $3.3\!\times\!10^{-4}$ & $2.7\!\times\!10^{-3}$ \\
 &  & 2 & $2.6\!\times\!10^{-5}$ & $2.2\!\times\!10^{-5}$ & 2.42 & $5.6\!\times\!10^{-5}$ & $3.4\!\times\!10^{-4}$ \\
J & LMG, $J{=}3/2$ & 0 & $1.9\!\times\!10^{-6}$ & $1.1\!\times\!10^{-6}$ & $1.3\!\times\!10^{-6}$ & $1.5\!\times\!10^{-6}$ & $2.2\!\times\!10^{-6}$ \\
 &  & 1 & $5.5\!\times\!10^{-6}$ & $1.1\!\times\!10^{-6}$ & $2.3\!\times\!10^{-6}$ & $4.5\!\times\!10^{-6}$ & $2.9\!\times\!10^{-6}$ \\
 &  & 2 & $6.2\!\times\!10^{-7}$ & $2.9\!\times\!10^{-6}$ & $3.3\!\times\!10^{-6}$ & $1.3\!\times\!10^{-6}$ & $6.5\!\times\!10^{-6}$ \\
\botrule
\end{longtable}


\section{$\mathbb{Z}_2$ gauge theory: penalty-strength sweep}\label{sec:supp_vsweep}

The main text's section on structural gauge invariance states that, on the penalty-enforced
hardware-efficient ansatz, Q-MAML is behind the best classical scheme at every penalty strength $V$
tested, refuting the explanation that a large penalty simply drowns the task signal.
Table~\ref{tab:supp_vsweep} gives the single-seed final energies behind that statement (reduced
4-task, 200-iteration budget, $V\in\{1,3,10,30\}$); Q-MAML ties Uniform only at the weakest penalty
and falls further behind as $V$ grows to $10$, recovering slightly by $V=30$.

\begin{table}[!htb]
\centering
\caption{$\mathbb{Z}_2$ gauge theory, penalty-enforced ansatz: final energy by penalty strength $V$ (single seed, reduced budget).}\label{tab:supp_vsweep}
\begin{tabular}{@{}lcccccc@{}}
\toprule
$V$ & Q-MAML & Zero & $\pi$ & Uniform & Gaussian & Q-MAML / best classical \\
\midrule
1 & 1.468 & 2.490 & 1.560 & 1.438 & 1.516 & 1.02 (Uniform) \\
3 & 2.683 & 3.335 & 2.280 & 1.655 & 2.042 & 1.62 (Uniform) \\
10 & 3.791 & 4.432 & 2.917 & 2.850 & 2.891 & 1.33 (Uniform) \\
30 & 4.157 & 5.498 & 3.729 & 3.657 & 4.165 & 1.14 (Uniform) \\
\botrule
\end{tabular}
\end{table}


\section{Lattice Schwinger model: the $L{=}4$ depth/pretraining-rate scan}\label{sec:supp_l4}

The main text's section on system size reports that depth, not pretraining convergence, decides the
$L{=}4$ (8-qubit) result. Table~\ref{tab:supp_l4} gives the single-seed grid this was read off:
pretraining learning rate 0.005 stalls (final energy far from the ground state at every depth-12
setting, visible as the shared Zero/Uniform/Gaussian columns not moving with epoch count), and at
depth 12 no learning-rate or epoch-count change brings Q-MAML ahead of $\pi$; only the depth-18
ansatz does, and by a wide margin.

\begin{table}[!htb]
\centering
\caption{Lattice Schwinger model, $L{=}4$: final relative error (\%) by depth and pretraining setting (single seed, reduced budget).}\label{tab:supp_l4}
\begin{tabular}{@{}lccccc@{}}
\toprule
Pretraining setting & Q-MAML & Zero & $\pi$ & Uniform & Gaussian \\
\midrule
Depth 12, 40 ep, lr 0.005 & 24.49 & 451 & 17.95 & 39.44 & 38.40 \\
Depth 12, 100 ep, lr 0.005 & 23.89 & 451 & 17.95 & 39.44 & 38.40 \\
Depth 12, 100 ep, lr 0.002 & 21.21 & 451 & 17.95 & 39.44 & 38.40 \\
Depth 18, 100 ep, lr 0.002 & 2.43 & 451 & 8.65 & 10.95 & 11.84 \\
\botrule
\end{tabular}
\end{table}


\section{Quark--gluon jet tagging}\label{sec:supp_qg}

A second classification dataset, quark--gluon jet tagging (six qubits, two seeds), is mentioned in
the main text's HEP classification section as directional evidence only. Table~\ref{tab:supp_qg}
gives the full comparison: Q-MAML's validation accuracy and meta-loss against the three classical
baselines used elsewhere in classification (the first-order variant was not run on this dataset).

\begin{table}[!htb]
\centering
\caption{Quark--gluon jet tagging, six qubits ($n=2$ seeds): final validation accuracy and meta-loss.}\label{tab:supp_qg}
\begin{tabular}{@{}lcc@{}}
\toprule
Scheme & Validation accuracy & Meta-loss \\
\midrule
Zero & 0.594\,$\pm$\,0.021 & 0.693\,$\pm$\,0.004 \\
Uniform & 0.594\,$\pm$\,0.021 & 0.701\,$\pm$\,0.007 \\
Gaussian & 0.599\,$\pm$\,0.036 & 0.689\,$\pm$\,0.000 \\
Q-MAML & 0.776\,$\pm$\,0.005 & 0.308\,$\pm$\,0.003 \\
\botrule
\end{tabular}
\end{table}


\section{Software and compute}\label{sec:supp_software}

Circuits are simulated with PennyLane; the Learner, encoders and optimizers use PyTorch (Adam
throughout, double precision for the Learner). All results are noiseless statevector simulation; no
quantum hardware was used. Exact ground energies and the symmetry checks in Table~1 (main text) use
NumPy/SciPy dense or sparse diagonalization, run once per Hamiltonian module before any Q-MAML run
(\texttt{verify\_construction()} in each \texttt{hamiltonians\_*.py}). Code, notebooks and saved
results are at the repository given in the main text's Data and code availability statement.

\end{document}
